% System section (Springer LNCS version). Describes only what is implemented in roadscribe/.
\section{RoadScribe}
\label{sec:system}

RoadScribe turns one recorded trip into a short narrative (Fig.~\ref{fig:arch}).
Three principles shape the design:
every number and place in the text must trace to a fact in a ledger;
the language model chooses and phrases content but never supplies a value;
and when its output fails verification, a deterministic template narrator replaces it.

\begin{figure}[t]
\centering
% Architecture figure (TikZ). Needs \usetikzlibrary{positioning,arrows.meta,fit}.
\begin{tikzpicture}[
  font=\scriptsize,
  box/.style={draw, rounded corners=2pt, align=center, minimum height=1.05cm, text width=2.05cm, inner sep=2pt},
  arr/.style={-{Stealth[length=2mm]}, thick},
  node distance=0.32cm and 0.32cm]
\node[box] (sens) {IMU 100~Hz\\GNSS 1~Hz};
\node[box, right=of sens] (l1) {\textbf{Layer 1}\\mount, ISO 2631\\felt / road value\\events};
\node[box, right=of l1] (l2) {\textbf{Layer 2}\\5~m grid\\PELT segments\\road units};
\node[box, right=of l2] (l3) {\textbf{Layer 3}\\claim ledger\\token-budgeted\\capsule};
\node[box, below=0.55cm of l3] (slm) {SLM\\plan + narrative};
\node[box, left=of slm] (ver) {verifier\\(claims vs.\ ledger)};
\node[box, left=of ver] (out) {narrative};
\node[box, below=0.55cm of l1, dashed] (cal) {calibration on\\shared roads:\\vehicle gain, $\hat\gamma$};
\draw[arr] (sens) -- (l1);
\draw[arr] (l1) -- (l2);
\draw[arr] (l2) -- (l3);
\draw[arr] (l3) -- (slm);
\draw[arr] (slm) -- (ver);
\draw[arr] (ver) -- node[above]{pass} (out);
\draw[arr, dashed] (cal) -- (l1);
\draw[arr] (l3.south west) ++(0.15,0) |- ++(-1.1,-0.25) -| node[pos=0.25, right, align=left]{template\\fallback} (ver.north);
\end{tikzpicture}

\caption{RoadScribe. Values flow from the sensors through Layers~1--3 into the claim ledger; the language model writes a plan and a narrative over ledger facts, and a verifier either accepts the narrative or falls back to the template narrator. The calibration block (dashed) is the remedy studied in Sect.~\ref{sec:ident}.}
\label{fig:arch}
\end{figure}

\subsection{Layer 1: vibration measures}
The IMU gives specific force $\mathbf f^b$ in the sensor frame at $f_s=100$~Hz and GNSS gives speed at 1~Hz.
The mount rotation $\mathbf R_{vb}$ is estimated without fixtures: the up axis is the geometric median of quasi-static gravity samples (Weiszfeld iteration with the Vardi--Zhang correction~\cite{vardi2000l1median}), and the remaining yaw is the closed-form orthogonal Procrustes fit~\cite{schonemann1966procrustes} of the levelled horizontal acceleration to GNSS longitudinal and centripetal acceleration.
The gravity-free vehicle-frame acceleration is $\mathbf a^v=\mathbf R_{vb}\mathbf f^b-g\mathbf e_z$.
ISO~2631-1~\cite{iso2631} weightings ($W_k$ vertical, $W_d$ horizontal) are applied exactly in the frequency domain, which is possible because narratives are written after the trip.
For a stretch $u$ the weighted RMS values give the \emph{felt} vibration $a_v=\sqrt{a_{wx}^2+a_{wy}^2+a_{wz}^2}$, and the original design defines a speed-decoupled \emph{road} value
\begin{equation}
\tilde a_{w}=a_{wz}\,(v_{\mathrm{ref}}/\bar v)^{\gamma},\qquad v_{\mathrm{ref}}=10~\mathrm{m\,s^{-1}},
\label{eq:dec}
\end{equation}
reported only for $\bar v\ge5~\mathrm{m\,s^{-1}}$.
The elasticity $\gamma$ follows from road spectra of the form $G_d(n)\propto n^{-w}$~\cite{iso8608,dodds1973roughness}: the RMS response of a fixed vehicle scales as $\bar v^{(w-1)/2}$, so $\gamma=(w-1)/2$.
Both values are mapped to the six comfort classes of ISO~2631-1 Annex~C with single cut-points $(0.315, 0.561, 0.894, 1.414, 2.236)~\mathrm{m\,s^{-2}}$.
Impulsive events are detected with the Teager--Kaiser energy operator~\cite{kaiser1990tkeo} on the 4--25~Hz vertical band against a robust running threshold, and labelled with their maneuver context from GNSS and the gyroscope.

\subsection{Layer 2: segmentation}
Features are resampled on a 5~m along-track grid and the trip is segmented by penalized change-point detection~\cite{killick2012pelt},
$\min_{\tau}\sum_i C(\mathbf y_{\tau_{i-1}+1:\tau_i})+\sum_i\beta(\tau_i)$,
with a Gaussian mean-shift cost.
The design lowers the penalty $\beta(\tau)$ near map transitions.
With a location-dependent penalty the textbook PELT pruning test is not exact; the correct test is $G(t)+C(\mathbf y_{t+1:s})\ge G(s)$ with $G(t)=F(t)+\beta(t)$, a direct generalization of PELT's argument~\cite{killick2012pelt}.
In 300 random series the textbook test returned a sub-optimal partition in 136, while the corrected test matched exhaustive optimal partitioning in all of them.
Map matching is not implemented in the version evaluated here, so on PVS the penalty is constant and the segmentation is plain PELT.
Short segments are merged, and each resulting road unit carries its Layer~1 measures and events.

\subsection{Layer 3: ledger, capsule and generation}
Units and events are compiled into a ledger of typed facts: severity classes (felt and road), measured values with tolerances, events with position and confidence, place anchors, comparisons, and \emph{missing} facts that record why a value is unavailable (e.g.\ speed too low to separate road from speed).
A capsule serializes the ledger within a token budget; units are chosen by an exact multiple-choice knapsack over salience.
The small language model (Qwen2.5-3B-Instruct~\cite{qwen25} with a QLoRA adapter~\cite{dettmers2023qlora}) receives the capsule and an audience (persona) and writes a plan over ledger identifiers followed by the narrative.
A closed-vocabulary verifier extracts every claim from the narrative and checks it against the ledger: numbers within tolerance, severity words within their comfort band, events matched to ledger events, no value stated for a missing field.
If the plan does not parse or any claim fails, the system returns the template narrator's text for the same capsule; this is a run-time-assurance pattern in the sense of the Simplex architecture~\cite{sha2001simplex}.
The adapter was fine-tuned on 2{,}696 synthetic capsule--narrative pairs whose narratives come from the template narrator (2{,}500) and from a frontier model writing in the verifier's lexicon (196).
